
\begin{tikzpicture}[scale=1.2]

% ============================================
% 1. PUNKT
% ============================================
\begin{scope}[shift={(-4,0)}]
    \fill (0,0) circle[radius=3pt];
\end{scope}

% ============================================
% 2. ODCINEK
% ============================================
\begin{scope}[shift={(-2,0)}]
    \draw[thick] (0,-1) -- (0,1);
    \fill (0,-1) circle[radius=2pt];
    \fill (0,1) circle[radius=2pt];
\end{scope}

% ============================================
% 3. TRÓJKĄT
% ============================================
\begin{scope}[shift={(0,-0.5)}]
     \fill[gray!20] (0,0) -- (2,0) -- (1,1.5) -- cycle;
    \draw[thick] (0,0) -- (2,0) -- (1,1.5) -- cycle;
    \fill (0,0) circle[radius=2pt];
    \fill (2,0) circle[radius=2pt];
    \fill (1,1.5) circle[radius=2pt];
\end{scope}

% ============================================
% 4. OSTROSŁUP TRÓJKĄTNY
% ============================================
\begin{scope}[shift={(4,-0.5)}]
    \tdplotsetmaincoords{70}{120}
    \begin{scope}[tdplot_main_coords, scale=0.8]
        
        % Wierzchołki
        \coordinate (A) at (0,0,0);
        \coordinate (B) at (2,0,0);
        \coordinate (C) at (0.5,1.5,0);
        \coordinate (S) at (0.8,0.6,2.5);
        
        \fill[gray!20] (S) -- (C) -- (B) -- cycle;
        
        % Krawędzie podstawy
        \draw[thick] (B)-- (C);
        \draw[dashed] (B) -- (A);
        \draw[dashed] (C) -- (A);
        
        % Krawędzie boczne
        \draw[thick] (S) -- (C);
        \draw[thick] (S) -- (B);
        \draw[dashed] (S) -- (A);
        
        % Wierzchołki
        \fill (A) circle[radius=2pt];
        \fill (B) circle[radius=2pt];
        \fill (C) circle[radius=2pt];
        \fill (S) circle[radius=2pt];
        
    \end{scope}
\end{scope}

\end{tikzpicture}